% fig11_multimodal_architecture.tex — 图11 多模态大模型标准架构
% 《深度学习与大语言模型：前沿技术全景报告》图示库
% 由 dl_llm_report.tex 抽取，经 % fig11_multimodal_architecture.tex — 图11 多模态大模型标准架构
% 《深度学习与大语言模型：前沿技术全景报告》图示库
% 由 dl_llm_report.tex 抽取，经 % fig11_multimodal_architecture.tex — 图11 多模态大模型标准架构
% 《深度学习与大语言模型：前沿技术全景报告》图示库
% 由 dl_llm_report.tex 抽取，经 % fig11_multimodal_architecture.tex — 图11 多模态大模型标准架构
% 《深度学习与大语言模型：前沿技术全景报告》图示库
% 由 dl_llm_report.tex 抽取，经 \input{figs/fig11_multimodal_architecture} 引用（caption/label 在主文件）
% 注意：本文件为片段，依赖主文件导言区的 tikz/pgfplots 宏包、颜色定义与 tikzset 样式，不可单独编译。

\begin{tikzpicture}
  % ---- 输入模态（左） ----
  \node[gry, minimum width=22mm] (img) at (0,1.7) {图像 / 文档};
  \node[gry, minimum width=22mm] (vid) at (0,0) {视频（帧序列）};
  \node[gry, minimum width=22mm] (aud) at (0,-1.7) {音频};
  % ---- 专用编码器 ----
  \node[blk, minimum width=26mm] (venc) at (3.5,1.1) {视觉编码器\\ ViT / ConvNeXt};
  \node[blk, minimum width=26mm] (aenc) at (3.5,-1.1) {音频编码器};
  % ---- 投影层 / LLM / 输出：阶段间隙 ≥0.7cm，避免贴挤 ----
  \node[orn, minimum width=26mm] (proj) at (6.9,0) {投影层\\ 对齐到词表空间};
  \node[grn, minimum width=28mm, minimum height=13mm] (llm) at (10.3,0) {LLM Backbone\\（文本 token 直连）};
  \node[orn, minimum width=28mm] (out) at (13.9,0) {输出\\ 文本 · 动作 · 工具调用};
  \node[gry, minimum width=22mm] (txt) at (6.9,2.9) {文本 token};
  % ---- 连线：编码器入口/投影层入口均分散锚点，箭头不再汇聚打结 ----
  \draw[arr] (img.east) -- ([yshift=2mm]venc.west);
  \draw[arr] (vid.east) -- ([yshift=-2mm]venc.west);
  \draw[arr] (aud.east) -- (aenc.west);
  \draw[arr] (venc.east) -- ([yshift=2mm]proj.west);
  \draw[arr] (aenc.east) -- ([yshift=-2mm]proj.west);
  \draw[arr] (txt) -- (proj);
  \draw[arr] (txt.east) to[bend left=22] (llm.70);
  \draw[arr] (proj)--(llm); \draw[arr] (llm)--(out);
  \node[lab, anchor=east] at (5.0,2.9) {拼接式（LLaVA 系）};
  \node[lab, anchor=west] at (10.3,2.1) {原生式（Gemini 系）};
  \node[lab] at (7.1,-3.1) {视觉 token 成本：一张 1024$^2$ 图 $\approx$ 数百至上千 token，是多模态请求成本的大头};
\end{tikzpicture}
 引用（caption/label 在主文件）
% 注意：本文件为片段，依赖主文件导言区的 tikz/pgfplots 宏包、颜色定义与 tikzset 样式，不可单独编译。

\begin{tikzpicture}
  % ---- 输入模态（左） ----
  \node[gry, minimum width=22mm] (img) at (0,1.7) {图像 / 文档};
  \node[gry, minimum width=22mm] (vid) at (0,0) {视频（帧序列）};
  \node[gry, minimum width=22mm] (aud) at (0,-1.7) {音频};
  % ---- 专用编码器 ----
  \node[blk, minimum width=26mm] (venc) at (3.5,1.1) {视觉编码器\\ ViT / ConvNeXt};
  \node[blk, minimum width=26mm] (aenc) at (3.5,-1.1) {音频编码器};
  % ---- 投影层 / LLM / 输出：阶段间隙 ≥0.7cm，避免贴挤 ----
  \node[orn, minimum width=26mm] (proj) at (6.9,0) {投影层\\ 对齐到词表空间};
  \node[grn, minimum width=28mm, minimum height=13mm] (llm) at (10.3,0) {LLM Backbone\\（文本 token 直连）};
  \node[orn, minimum width=28mm] (out) at (13.9,0) {输出\\ 文本 · 动作 · 工具调用};
  \node[gry, minimum width=22mm] (txt) at (6.9,2.9) {文本 token};
  % ---- 连线：编码器入口/投影层入口均分散锚点，箭头不再汇聚打结 ----
  \draw[arr] (img.east) -- ([yshift=2mm]venc.west);
  \draw[arr] (vid.east) -- ([yshift=-2mm]venc.west);
  \draw[arr] (aud.east) -- (aenc.west);
  \draw[arr] (venc.east) -- ([yshift=2mm]proj.west);
  \draw[arr] (aenc.east) -- ([yshift=-2mm]proj.west);
  \draw[arr] (txt) -- (proj);
  \draw[arr] (txt.east) to[bend left=22] (llm.70);
  \draw[arr] (proj)--(llm); \draw[arr] (llm)--(out);
  \node[lab, anchor=east] at (5.0,2.9) {拼接式（LLaVA 系）};
  \node[lab, anchor=west] at (10.3,2.1) {原生式（Gemini 系）};
  \node[lab] at (7.1,-3.1) {视觉 token 成本：一张 1024$^2$ 图 $\approx$ 数百至上千 token，是多模态请求成本的大头};
\end{tikzpicture}
 引用（caption/label 在主文件）
% 注意：本文件为片段，依赖主文件导言区的 tikz/pgfplots 宏包、颜色定义与 tikzset 样式，不可单独编译。

\begin{tikzpicture}
  % ---- 输入模态（左） ----
  \node[gry, minimum width=22mm] (img) at (0,1.7) {图像 / 文档};
  \node[gry, minimum width=22mm] (vid) at (0,0) {视频（帧序列）};
  \node[gry, minimum width=22mm] (aud) at (0,-1.7) {音频};
  % ---- 专用编码器 ----
  \node[blk, minimum width=26mm] (venc) at (3.5,1.1) {视觉编码器\\ ViT / ConvNeXt};
  \node[blk, minimum width=26mm] (aenc) at (3.5,-1.1) {音频编码器};
  % ---- 投影层 / LLM / 输出：阶段间隙 ≥0.7cm，避免贴挤 ----
  \node[orn, minimum width=26mm] (proj) at (6.9,0) {投影层\\ 对齐到词表空间};
  \node[grn, minimum width=28mm, minimum height=13mm] (llm) at (10.3,0) {LLM Backbone\\（文本 token 直连）};
  \node[orn, minimum width=28mm] (out) at (13.9,0) {输出\\ 文本 · 动作 · 工具调用};
  \node[gry, minimum width=22mm] (txt) at (6.9,2.9) {文本 token};
  % ---- 连线：编码器入口/投影层入口均分散锚点，箭头不再汇聚打结 ----
  \draw[arr] (img.east) -- ([yshift=2mm]venc.west);
  \draw[arr] (vid.east) -- ([yshift=-2mm]venc.west);
  \draw[arr] (aud.east) -- (aenc.west);
  \draw[arr] (venc.east) -- ([yshift=2mm]proj.west);
  \draw[arr] (aenc.east) -- ([yshift=-2mm]proj.west);
  \draw[arr] (txt) -- (proj);
  \draw[arr] (txt.east) to[bend left=22] (llm.70);
  \draw[arr] (proj)--(llm); \draw[arr] (llm)--(out);
  \node[lab, anchor=east] at (5.0,2.9) {拼接式（LLaVA 系）};
  \node[lab, anchor=west] at (10.3,2.1) {原生式（Gemini 系）};
  \node[lab] at (7.1,-3.1) {视觉 token 成本：一张 1024$^2$ 图 $\approx$ 数百至上千 token，是多模态请求成本的大头};
\end{tikzpicture}
 引用（caption/label 在主文件）
% 注意：本文件为片段，依赖主文件导言区的 tikz/pgfplots 宏包、颜色定义与 tikzset 样式，不可单独编译。

\begin{tikzpicture}
  % ---- 输入模态（左） ----
  \node[gry, minimum width=22mm] (img) at (0,1.7) {图像 / 文档};
  \node[gry, minimum width=22mm] (vid) at (0,0) {视频（帧序列）};
  \node[gry, minimum width=22mm] (aud) at (0,-1.7) {音频};
  % ---- 专用编码器 ----
  \node[blk, minimum width=26mm] (venc) at (3.5,1.1) {视觉编码器\\ ViT / ConvNeXt};
  \node[blk, minimum width=26mm] (aenc) at (3.5,-1.1) {音频编码器};
  % ---- 投影层 / LLM / 输出：阶段间隙 ≥0.7cm，避免贴挤 ----
  \node[orn, minimum width=26mm] (proj) at (6.9,0) {投影层\\ 对齐到词表空间};
  \node[grn, minimum width=28mm, minimum height=13mm] (llm) at (10.3,0) {LLM Backbone\\（文本 token 直连）};
  \node[orn, minimum width=28mm] (out) at (13.9,0) {输出\\ 文本 · 动作 · 工具调用};
  \node[gry, minimum width=22mm] (txt) at (6.9,2.9) {文本 token};
  % ---- 连线：编码器入口/投影层入口均分散锚点，箭头不再汇聚打结 ----
  \draw[arr] (img.east) -- ([yshift=2mm]venc.west);
  \draw[arr] (vid.east) -- ([yshift=-2mm]venc.west);
  \draw[arr] (aud.east) -- (aenc.west);
  \draw[arr] (venc.east) -- ([yshift=2mm]proj.west);
  \draw[arr] (aenc.east) -- ([yshift=-2mm]proj.west);
  \draw[arr] (txt) -- (proj);
  \draw[arr] (txt.east) to[bend left=22] (llm.70);
  \draw[arr] (proj)--(llm); \draw[arr] (llm)--(out);
  \node[lab, anchor=east] at (5.0,2.9) {拼接式（LLaVA 系）};
  \node[lab, anchor=west] at (10.3,2.1) {原生式（Gemini 系）};
  \node[lab] at (7.1,-3.1) {视觉 token 成本：一张 1024$^2$ 图 $\approx$ 数百至上千 token，是多模态请求成本的大头};
\end{tikzpicture}
